\origpage{310}{ROBERTSON AND SEYMOUR}
\noindent that $G$ is a subgraph of a ``chordal'' graph with all cliques of size at most $w + 1$.)

Then, for example, trees and forests have tree-width 1 (or 0 in degenerate cases) and it can be shown that series-parallel graphs have tree-width $\le 2$. For $n \ge 1$, the complete graph $K_n$ has tree-width $n - 1$, while the $n \times n$ grid (this has $n^2$ vertices---it is the adjacency graph of the chessboard with $n^2$ squares), for $n \ge 2$, has tree-width $n$.

A graph $G$ is \textit{contractible} to $H$ when $H$ can be obtained from $G$ by contracting edges. $H$ is a \textit{minor} of $G$ if $G$ has a subgraph contractible to $H$. Other papers in this series [9, 10, 12, 14] establish a connection between the absence of a fixed graph $H$ as a minor of an arbitrary graph $G$ and the presence of a fixed type of structure on $G$. For example, in [12] it is shown that \textit{if} $H$ \textit{is a planar graph there is a number} $w$, \textit{depending only on} $H$, \textit{such that any graph} $G$ \textit{with no minor isomorphic to} $H$ \textit{has tree-width} $\le w$. Such structure can be exploited to obtain, by developments of standard methods, results which have defied solution for graphs in general. Thus, in [11] it is shown that \textit{any set of graphs with bounded tree-width, no member of which is isomorphic to a minor of another member, must be finite}. This was proved by Kruskal [5] for graphs of tree-width $\le 1$, under the stronger inclusion relation of topological containment.

The main object of this paper is to use the result from [12] to obtain an algorithm for testing minor inclusion of a fixed planar graph $H$. The \textit{size} of a graph $G$ is $|V(G)| + |E(G)|$. If the running time of a graph algorithm can be bounded above by a polynomial in the size of the input graph, it is called a \textit{polynomial algorithm}. Our main results are as follows.

\stmt
(1.1) \textit{For any fixed planar graph} $H$, \textit{there is a polynomial algorithm to decide if the input graph has a minor isomorphic to} $H$.

\stmt
(1.2) \textit{For any fixed integer} $w$, \textit{there is a polynomial algorithm to decide if the input graph has tree-width} $\le w$.

\stmt
(1.3) \textit{For any fixed integers} $k, w$ \textit{there is a polynomial algorithm to solve the} DISJOINT CONNECTING PATHS \textit{problem (stated below) for} $k$ \textit{paths, in graphs of tree-width} $\le w$.

Our proofs for (1.1) and (1.3) are constructive (we find the algorithms), but for (1.2) the proof is nonconstructive. We should stress that these make no claim to be practical algorithms, for the exponent of the polynomial is very large, even for quite small $H, k, w$. Our concern is more theoretical, with the \textit{NP}-completeness (or lack of it) of problems of this type. (We shall assume familiarity with this concept---see [2].)

A \textit{termination} $\mathscr{A}$ in a graph $G$ is a set $\mathscr{A} = \{(s_i, t_i) : 1 \le i \le k\}$ of pairs of vertices of $G$, such that $s_1, t_1, s_2, t_2, \dots, s_k, t_k$ are all distinct. The \textit{support} $S(\mathscr{A})$ of $\mathscr{A}$ is $\{s_1, t_1, \dots, s_k, t_k\}$. The pair $(G, \mathscr{A})$ is a \textit{terminated graph} and it has \textit{index} $k$. The terminated graph is \textit{feasible} if there are paths
